\section{The local-time potential}\label{sec:local}
We represent occupation times by a potential whose spherical averages
measure the part of the range outside a ball.

Let $P$ be the simple random walk transition operator on $\Z^d$.
For $d=2$, let $b$ be the potential kernel normalized by $b(0)=0$.
For $d\geq3$, let $b=-G$, where
$G(x)=\sum_{j\geq0}P^j(0,x)$ is the Green function.
In both cases $(P-I)b=\ind_{\{0\}}$.
The lattice potential estimates of
\citet[Theorem~4.3.1, Corollary~4.3.3, and Theorem~4.4.4]{LawlerLimic2010}
give, as $|x|\to\infty$,
\begin{equation}\label{eq:kernel-asymptotics}
b(x)=\begin{cases}
\dfrac2\pi\log|x|+\kappa+O(|x|^{-2}),&d=2,\\[3pt]
-\dfrac{2}{(d-2)\omega_d}|x|^{2-d}+O_d(|x|^{-d}),&d\geq3
\end{cases}\, ,
\end{equation}
where $\kappa$ is the planar potential-kernel constant.
For the central difference $D_i b(x)=(b(x+e_i)-b(x-e_i))/2$,
\begin{equation}\label{eq:gradient}
Db(x)=\frac2{\omega_d}\frac{x}{|x|^d}+O_d(|x|^{-d}) ,
\qquad |b(x+e)-b(x)|\leq C_d(1+|x|)^{1-d}\, .
\end{equation}
The first estimate concerns $x\ne0$; the second concerns every lattice
site and every unit coordinate vector with either sign.

We use the following consequence of \citet[Theorem~1.6]{Freedman1975}.
If the increments of a real martingale $Z$ have absolute value at most
$B$, then, for deterministic $v,t>0$,
\begin{equation}\label{eq:freedman}
\P\bigl(|Z_n-Z_0|\geq t,\ \langle Z\rangle_n\leq v\bigr)
\leq2\exp\left(-\frac{t^2}{2(v+Bt)}\right)\, .
\end{equation}
We apply this inequality at dyadic variance bounds and take polynomial
union bounds over deterministic targets and time intervals.
The resulting estimates can then be used at points and intervals
selected from the path.

For $m\geq0$, define the error scale
\begin{equation}\label{eq:error-scale}
e_n(m)=\begin{cases}
\sqrt m L+L,&d=2,\\
\sqrt{mL}+L,&d\geq3
\end{cases}
\qquad
\lambda_n=\begin{cases}L^2,&d=2,\\L,&d\geq3
\end{cases}\, .
\end{equation}
For integers $0\leq s<t\leq n$, let
\begin{equation}\label{eq:interval-def}
\ell_{s,t}(x)=\sum_{s\leq j<t}\ind_{\{X_j=x\}} ,\qquad
M_{s,t}=\max_x\ell_{s,t}(x) ,\qquad
k_{s,t}=\sum_{s\leq j<t}I_j\, .
\end{equation}

\begin{lemma}\label{lem:local}
Let $d\geq2$ be an integer, let $0<\eps<1/d$, and let $(X_j)_{j\geq0}$ be
centrally excited random walk started at the origin, with first-departure mean $-\eps u_x$
and probabilities $q_x(\pm e_i)=1/(2d)\mp\eps u_x^i/2$,
with simple symmetric steps on later departures and at the origin.
For every $p>0$ there is $C=C(d,\eps,p)$ such that, for every integer
$n\geq2$, with probability at least $1-Cn^{-p}$, the following estimates
hold simultaneously:
\begin{align}
M_n&\leq C_d\eps R_n^{1/d}+C\lambda_n\, ,\label{eq:M}\\
M_{s,t}&\leq C_d\eps k_{s,t}^{1/d}+C\lambda_n
\qquad(0\leq s<t\leq n)\, ,\label{eq:interval}\\
|\widetilde\ell_n(y)-U_{D_n}(y)|&\leq C e_n(M_n)
\qquad(y\in\R^d,\ |y|\leq2n)\, .\label{eq:approx}
\end{align}
Here $U_D(y)=(2\eps/\omega_d)\int_D u_v\cdot(v-y)|v-y|^{-d}\dd v$,
$L=\log(n+2)$, and $\lambda_n,e_n$ have the values displayed above.
\end{lemma}

\begin{proof}
Dynkin's formula for $b(X_j-y)$ gives a martingale $\mathcal M^y$ with
\begin{equation}\label{eq:dynkin}
\ell_n(y)=b(X_n-y)-b(-y)
+\eps\sum_{x\in A_n}u_x\cdot Db(x-y)-\mathcal M_n^y\, .
\end{equation}
The endpoint difference is $O(L)$ for $d=2$ and $O_d(1)$ for $d\geq3$,
uniformly for lattice targets $|y|\leq3n$.
The increments of $\mathcal M^y$ are bounded, and
\begin{equation}\label{eq:bracket}
\langle\mathcal M^y\rangle_n\leq C_d B_y ,\qquad
B_y=\sum_x\ell_n(x)(1+|x-y|)^{2-2d}\, .
\end{equation}
The sum of the weights over $|x-y|\leq4n$ is $O(L)$ for $d=2$
and $O_d(1)$ for $d\geq3$. Applying \eqref{eq:freedman} at dyadic
bounds for $M_{s,t}\leq n$, and taking a union bound over the
$O_d(n^{d+2})$ interval--target pairs, gives
\begin{equation}\label{eq:interval-mart}
|\mathcal M_t^y-\mathcal M_s^y|\leq C e_n(M_{s,t})\, .
\end{equation}
For every set $E$ of $k$ lattice sites, filling concentric shells gives
\begin{equation}\label{eq:packing-lattice}
\sup_y\sum_{x\in E}(1+|x-y|)^{1-d}\leq C_d k^{1/d}\, .
\end{equation}
Subtract \eqref{eq:dynkin} at $s$ and $t$ and maximize over $y$.
The source sum contains $k_{s,t}$ distinct sites. Young's inequality
absorbs the square root of $M_{s,t}$ and proves \eqref{eq:interval}.
The choice $s=0,t=n$ proves \eqref{eq:M}.

To pass from the lattice sum to $U_{D_n}$, first replace $Db(x-y)$
by $(2/\omega_d)(x-y)/|x-y|^d$. The errors sum to
$C\sum_{|x|\leq n}(1+|x-y|)^{-d}\leq CL$.
Replacing each kernel value by its cell integral has the same bound:
the distant-cell error is $O_d((1+|x-y|)^{-d})$, and the kernel
$|v-y|^{1-d}$ is locally integrable near the remaining cells.
Since $|u_x-u_v|\leq C_d/(1+|v|)$ on $C_x$, replacing the direction
costs at most
\begin{equation}\label{eq:direction-error}
C_d\int_{B(0,n+\sqrt d)}
\frac{\dd v}{(1+|v|)|v-y|^{d-1}}\leq C_dL\, .
\end{equation}
For this bound, remove $|v-y|<1$ and use H\"older's inequality with
exponents $d$ and $d/(d-1)$; both norm powers are logarithmic.
Moving $y$ inside its cell costs another $O_d(L)$ by the same
near/far decomposition. Equation~\eqref{eq:approx} follows.
\end{proof}

We will retain the variance at individual targets. Since $B_y\leq n$,
dyadic bounds for the bracket in \eqref{eq:bracket} give, on an event
with the same probability bound,
\begin{equation}\label{eq:localmart}
|\mathcal M_n^y|\leq C(\sqrt{B_yL}+L)
\qquad(y\in\Z^d,\ |y|\leq3n)\, .
\end{equation}
The deterministic replacements in the preceding proof also give
\begin{equation}\label{eq:pointwise}
\ell_n(y)=U_{D_n}(y)+\rho_n(y)-\mathcal M_n^y ,\qquad
|\rho_n(y)|\leq CL\, .
\end{equation}
These formulas apply simultaneously to the same deterministic targets.
For real points $y,z$, we also use the cell modulus
\begin{equation}\label{eq:cellmodulus}
|U_{D_n}(y)-U_{D_n}(z)|\leq CL
\qquad(|y-z|\leq\sqrt d,\ |y|,|z|\leq2n)\, .
\end{equation}

\subsection{Spherical averages}
The radial average removes the angular dependence of the potential.
Let $\sigma_d=d\omega_d$ be the surface area of the unit sphere.
For a bounded measurable set $D\subset\R^d$, write
\begin{equation}\label{eq:Fdef}
f(r)=\frac1{\sigma_d}\int_{S^{d-1}}\ind_D(r\theta)\dd S(\theta) ,
\qquad F(s)=\int_s^\infty f(r)\dd r\, .
\end{equation}
Then $F$ is nonincreasing and absolutely continuous, and
\begin{equation}\label{eq:Fmass}
F(s)=\frac1{\sigma_d}\int_{D\cap\{|v|>s\}}|v|^{1-d}\dd v
\leq\frac{|D|}{\sigma_ds^{d-1}}
\qquad(s>0)\, .
\end{equation}

\begin{lemma}\label{lem:geometry}
Let $d\geq2$ be an integer, let $\eps>0$, and let $D\subset\R^d$
be a bounded measurable set of volume $R$. Let
$U_D(y)=(2\eps/\omega_d)\int_Du_v\cdot(v-y)|v-y|^{-d}\dd v$,
and let $F(s)=\sigma_d^{-1}\int_{D\cap\{|v|>s\}}|v|^{1-d}\dd v$.
Then there is a constant $C_d>0$, depending only on $d$, such that
for every $y,z\in\R^d$ and every $s>0$,
\begin{align}
\|U_D\|_\infty&\leq C_d\eps R^{1/d}\, ,\label{eq:potential-bound}\\
|U_D(y)-U_D(z)|&\leq C_d\eps R^{1/(2d)}|y-z|^{1/2}\, ,\label{eq:holder}\\
\frac1{\sigma_d}\int_{S^{d-1}}U_D(s\theta)\dd S(\theta)
&=2d\eps F(s)\, .\label{eq:newton}
\end{align}
For every $b>0$, the ball potential is
\begin{equation}\label{eq:ballpotential}
U_{B(0,b)}(y)=2d\eps(b-|y|)_+\, .
\end{equation}
\end{lemma}

\begin{proof}
Integrating $|v-y|^{1-d}$ over a ball of volume $R$ centered at $y$ proves
\eqref{eq:potential-bound}. For $K(v)=v/|v|^d$ and
$q=2d/(2d-1)$, scaling gives
\begin{equation}\label{eq:kernel-modulus}
\|K(\cdot-y)-K(\cdot-z)\|_{L^q(\R^d)}=C_d|y-z|^{1/2}\, .
\end{equation}
The norm is finite: the singularities have order $d-1$, and the
difference at infinity has order $d$. H\"older's inequality with
the $L^{2d}$ norm of $\ind_Du$ proves \eqref{eq:holder}.

The spherical mean of $\log|v-s\theta|$ in dimension two is
$\log\max(|v|,s)$. For $d\geq3$, the spherical mean of
$|v-s\theta|^{2-d}$ is $\max(|v|,s)^{2-d}$.
These identities follow from the mean-value property and radial
harmonicity inside and outside the source sphere.
Differentiating with respect to $|v|$, away from $|v|=s$, gives
\begin{equation}\label{eq:kernel-average}
\frac1{\sigma_d}\int_{S^{d-1}}
u_v\cdot\frac{v-s\theta}{|v-s\theta|^d}\dd S(\theta)
=|v|^{1-d}\ind_{\{|v|>s\}}\, .
\end{equation}
Fubini's theorem applies because
$\sup_y\int_D|v-y|^{1-d}\dd v<\infty$.
Integration over $D$ proves \eqref{eq:newton}; continuity covers
every radius. For a ball, $U_D$ is radial, and
$F(s)=(b-s)_+$. This proves \eqref{eq:ballpotential}.
\end{proof}
